%------------------------------------------------------------------------------%
\documentclass[fleqn,utf8,aspectratio=169,12pt]{beamer}

\usepackage{gaal}

\title{Equação do Plano}

%------------------------------------------------------------------------------%
\begin{document}

\addtitle

%------------------------------------------------------------------------------%
\section{Equação Vetorial do Plano}

%------------------------------------------------------------------------------%
\begin{frame}{Plano Determinado por um Ponto e Dois Vetores}
  \onslide<+->{Um plano pode ser determinado por}
  \begin{itemize}[<+->]
      \item um ponto $A$ pertencente ao plano
      \item dois vetores não nulos e não paralelos $\vec{v}$ e $\vec{w}$
  \end{itemize}
\end{frame}

%------------------------------------------------------------------------------%
\framedgraphic{Plano Determinado por um Ponto e Dois Vetores}{F-03-Equacao_do_plano-1}{}
\framedgraphic{Plano Determinado por um Ponto e Dois Vetores}{F-03-Equacao_do_plano-2}{}
\framedgraphic{Plano Determinado por um Ponto e Dois Vetores}{F-03-Equacao_do_plano-3}{}
\framedgraphic{Plano Determinado por um Ponto e Dois Vetores}{F-03-Equacao_do_plano-4}{}

%------------------------------------------------------------------------------%
\begin{frame}{Plano Determinado por um Ponto e Dois Vetores}
  \onslide<+->{Se $P$ é um ponto qualquer do plano,}

  \onslide<+->{então o vetor $\overrightarrow{AP}$}

  \onslide<+->{pode ser escrito como combinação linear de $\vec{v}$ e $\vec{w}$}
  \begin{align*}
    \onslide<+->{\overrightarrow{AP} & = s\vec{v} + t\vec{w} \\}
    \onslide<+->{P - A & = s\vec{v} + t\vec{w} \\}
    \onslide<+->{P & = A + s\vec{v} + t\vec{w} }
  \end{align*}
  \onslide<+->{com $s$, $t\in\R$}
\end{frame}

%------------------------------------------------------------------------------%
\begin{frame}{Equação Vetorial do Plano}
  Se
  \(
    \pause
    A = \vetortri{x_0}{y_0}{z_0}
    \pause\qquad
    \vec{v} = \vetortri{a}{b}{c}
    \pause\qquad
    \vec{w} = \vetortri{d}{e}{f}
  \)

  \pause
  \smallskip
  a \emph{Equação Vetorial do Plano} é
  \[
    \pause
    \vec{r} = A + s\vec{v} + t\vec{w}
  \]
  \[
    \pause
    \vetortri{x}{y}{z}
    = \vetortri{x_0}{y_0}{z_0}
    + s\vetortri{a}{b}{c}
    + t\vetortri{d}{e}{f}
    \qquad\qquad
    s, \, t \in \R
  \]
\end{frame}

%------------------------------------------------------------------------------%
\begin{frame}{Equações Paramétricas do Plano}
  Reescrevendo
  \[
    \vetortri{x}{y}{z}
    = \vetortri{x_0}{y_0}{z_0}
    + s\vetortri{a}{b}{c}
    + t\vetortri{d}{e}{f}
  \]
  \pause
  compomente por componente,
  \pause
  obtemos as \emph{Equações Paramétricas do Plano}
  \[
    \pause
    \begin{cases}
      x = x_0 + as + dt \\
      y = y_0 + bs + et \\
      z = z_0 + cs + ft
    \end{cases}
  \]
\end{frame}

%------------------------------------------------------------------------------%
%------------------------------------------------------------------------------%
\begin{question}
  Determine uma equação vetorial e as equações paramétricas do plano que passa pelo ponto
  \[
    A = (1, 2, -1)
  \]
  e contém os vetores
  \[
    \vec{v} = (2, -1, 3)
  \]
  \[
    \vec{w} = (1, 2, 1)
  \]
\end{question}

%------------------------------------------------------------------------------%
\begin{resolution}{Solução}
  A equação vetorial do plano é
  \[
    \vec{r} = A + s\vec{v} + t\vec{w}
  \]
  \pause
  \[
    \vetortri{x}{y}{z}
    = \vetortri{1}{2}{-1}
    + s\vetortri{2}{-1}{3}
    + t\vetortri{1}{2}{1}
    \qquad\qquad
    s, \, t \in \R
  \]
  \pause
  Separando as coordenadas
  \[
    \begin{cases}
      x =  1 + 2s +  t  \\
      y =  2 -  s + 2t  \\
      z = -1 + 3s + t
    \end{cases}
  \]
\end{resolution}

%------------------------------------------------------------------------------%
\endinput


%------------------------------------------------------------------------------%
\section{Equação Geral do Plano}

%------------------------------------------------------------------------------%
\begin{frame}{Vetor Normal ao Plano}
  Qualquer vetor ortogonal ao plano é um vetor \emph{normal ao plano}

  \pause
  \bigskip
  Um vetor normal a um plano é ortogonal a todos os vetores contidos no plano

  \pause
  \bigskip
  Dados vetores $\vec{v}$ e $\vec{w}$ não nulos e não paralelos no plano

  \pause
  \bigskip
  \(\vec{v} \times \vec{w}\) é um vetor normal ao plano
\end{frame}

%------------------------------------------------------------------------------%
\begin{frame}{Equação Geral do Plano}
  Considere um plano que contém o \emph{ponto}
  \[
    A = (x_0, y_0, z_0)
  \]
  \pause
  e possui \emph{vetor normal}
  \[
    \vec{n} = \vetortri{a}{b}{c}
  \]
  \pause
  Se $P = (x, y, z)$ pertence ao plano
  \pause
  \(
    \overrightarrow{AP}
  \)
  é perpendicular a $\vec{n}$
\end{frame}

%------------------------------------------------------------------------------%
\begin{frame}{Equação Geral do Plano}
\begin{align*}
  \onslide<+->{\vec{n} \cdot \overrightarrow{AP}     & = 0 \\}
  \onslide<+->{\vetortri{a}{b}{c} \cdot \vetortri{x-x_0}{y-y_0}{z-z_0} & = 0 \\}
  \onslide<+->{a(x-x_0) + b(y-y_0) + c(z-z_0)        & = 0 \\}
  \onslide<+->{ax + by + cz + (-ax_0 - by_0 - cz_0)  & = 0 \\[3mm]}
  \onslide<+->{ax + by + cz + d                      & = 0}
\end{align*}
\end{frame}

%------------------------------------------------------------------------------%
%------------------------------------------------------------------------------%
\begin{question}
  Determine a equação geral do plano que passa pelo ponto
  \[
    A = (1, 2, -1)
  \]
  e contém os vetores
  \[
    \vec{v} = (2, -1, 3)
  \]
  \[
    \vec{w} = (1, 2, 1)
  \]
\end{question}

%------------------------------------------------------------------------------%
\begin{resolution}{Vetor Normal}
  \onslide<+->{Um vetor normal é}
  \begin{align*}
    \onslide<+->{\vec{n}}
    \onslide<+->{& = \vec{v}\times\vec{w} \\}
    \onslide<+->{& = \begin{vmatrix}
          \vec{i} & \vec{j} & \vec{k} \\
          2       & -1      & 3       \\
          1       & 2       & 1
        \end{vmatrix}}
  \\
  \onslide<+->{& = (-1-6, 3-2, 4+1) \\}
  \onslide<+->{& = (-7, 1, 5)}
  \end{align*}
\end{resolution}

%------------------------------------------------------------------------------%
\begin{resolution}{Solução}
  \onslide<+->{A equação do plano é}
  \begin{align*}
    \onslide<+->{\vec{n} \cdot \overrightarrow{AP} & = 0 \\}
    \onslide<+->{-7(x-1)+(y-2)+5(z+1)              & = 0 \\}
    \onslide<+->{-7x+7+y-2+5z+5                    & = 0 \\}
    \onslide<+->{-7x+y+5z+10                       & = 0}
  \end{align*}
\end{resolution}

%------------------------------------------------------------------------------%
\endinput


%------------------------------------------------------------------------------%
\begin{question}
  Determine a equação geral do plano que passa pelos pontos
  \[
    A = (1, 0, 2)
  \]
  \[
    B = (3, 1, 1)
  \]
  \[
    C = (0, 2, 3)
  \]
\end{question}

%------------------------------------------------------------------------------%
\begin{resolution}{Solução}
  Construímos dois vetores contidos no plano

  \pause
  \[
    \vec{v}
    \pause
    = \overrightarrow{AB}
    \pause
    = B - A
    \pause
    = (3, 1, 1) - (1, 0, 2)
    \pause
    = (2, 1, -1)
  \]

  \pause
  \[
    \vec{w}
    \pause
    = \overrightarrow{AC}
    \pause
    = C - A
    \pause
    = (0, 2, 3) - (1, 0, 2)
    \pause
    = (-1, 2, 1)
  \]
\end{resolution}

%------------------------------------------------------------------------------%
\begin{resolution}{Solução}
  \onslide<+->{Um vetor normal é}
  \begin{align*}
    \onslide<+->{\vec{n}}
    \onslide<+->{& = \vec{v} \times \vec{w} \\}
    \onslide<+->{& = \begin{vmatrix}
          \vec{i} & \vec{j} & \vec{k} \\
          2       & 1       & -1      \\
          -1      & 2       & 1
        \end{vmatrix}}
    \\
    \onslide<+->{& = (3, -1 ,5)}
  \end{align*}
\end{resolution}

%------------------------------------------------------------------------------%
\begin{resolution}{Solução}
  \onslide<+->{Usando o ponto $A=(1,0,2)$}
  \begin{align*}
    \onslide<+->{\vec{n} \cdot \overrightarrow{AP} & = 0 \\}
    \onslide<+->{(3, -1 ,5) \cdot (x-1, y-0, z-2)  & = 0 \\}
    \onslide<+->{3(x-1)-y+5(z-2)                   & = 0 \\}
    \onslide<+->{3x-3-y+5z-10                      & = 0 \\}
    \onslide<+->{3x-y+5z-13                        & = 0}
  \end{align*}
\end{resolution}

%------------------------------------------------------------------------------%
\endinput

%------------------------------------------------------------------------------%
\begin{question}
  Determine uma parametrização para o plano
  \[
    2x - y + z - 3 = 0
  \]
\end{question}

%------------------------------------------------------------------------------%
\begin{resolution}{Solução}
\begin{columns}[t]
\column{0.5\textwidth}
  Temos
  \[
    2x-y+z-3=0
  \]
  \pause
  Isolamos uma coordenada
  \[
    y=2x+z-3
  \]
  \pause
  Escolhemos duas variáveis livres
  \[
    x = s
  \]
  \[
    z = t
  \]
\column{0.5\textwidth}
  \pause
  Equação parametrica do plano
  \[
    \begin{cases}
      x = s \\
      y = -3 +2s + t \\
      z = t
    \end{cases}
  \]
  com $s$, $t\in\mathbb{R}$
\end{columns}
\end{resolution}

%------------------------------------------------------------------------------%
\endinput

%------------------------------------------------------------------------------%
\begin{question}
  Determine se os pontos
  \(
    P = (1, 2, 0)
  \)
  e
  \(
    Q = (1, 1, -1)
  \)
  pertencem ao plano
  \[
  \gamma\colon 2x - 3y + z + 4 = 0
  \]
\end{question}

%------------------------------------------------------------------------------%
\begin{resolution}{Solução}
  \onslide<+->{Substituímos as coordenadas de $P = (1, 2, 0)$ na equação}
  \begin{align*}
    \onslide<+->{2x - 3y + z + 4}
    \onslide<+->{& = 2(1) - 3(2) + 0 + 4 \\}
    \onslide<+->{& = 2 - 6 + 4 \\}
    \onslide<+->{& = 0}
  \end{align*}
  \onslide<+->{A equação do plano é satisfeita}
  \onslide<+->{e $P\in\gamma$}
\end{resolution}

%------------------------------------------------------------------------------%
\begin{resolution}{Solução}
  \onslide<+->{Substituímos as coordenadas de $Q = (1, 1, -1)$ na equação}
  \begin{align*}
    \onslide<+->{2x - 3y + z + 4}
    \onslide<+->{& = 2(1) - 3(1) + (-1) + 4 \\}
    \onslide<+->{& = 2 - 3 - 1 + 4 \\}
    \onslide<+->{& = 2 \\}
    \onslide<+->{& \neq 0}
  \end{align*}
  \onslide<+->{A equação do plano não é satisfeita}
  \onslide<+->{e $Q\notin\gamma$}
\end{resolution}

%------------------------------------------------------------------------------%
\endinput


%%------------------------------------------------------------------------------%
%\section{Lista Mínima}
%
%%------------------------------------------------------------------------------%
%\begin{frame}{Lista Mínima}
%  \begin{enumerate}
%    \item
%  \end{enumerate}
%
%  \vfill
%  Atenção: A prova é baseada no livro, não nas apresentações
%\end{frame}

%------------------------------------------------------------------------------%
\end{document}
%------------------------------------------------------------------------------%
